% =====================================================================
%  Extended captions for the six UTL validation panels.
%  Included via \input from the main paper inside a figure float.
% =====================================================================

\newcommand{\captionUTLpanelone}{%
\textbf{Calibration of the Universal Transport Law.}
(a) Predicted versus measured $TP^{-1}$ across $20$ randomly drawn synthetic
kernels with $N \in \{2,\ldots,5\}$ decision classes. Each point is a single
trial; the dashed red line is the identity $y=x$. The Monte Carlo estimator
recovers $TP^{-1} = N^{-2} \sum_{i,j} \tau^{(ij)} g^{(ij)}$ to within
floating-point tolerance ($\max$ rel.\ err.\ $<10^{-2}$ at $n_\mathrm{dec} = 2 \times 10^4$).
(b) Per-trial residual; the tolerance threshold of $0.10$ is shown for
reference, but every trial lies far below it.
(c) Histogram of relative errors across all trials.
(d) Three-dimensional view of the integrand surface
$\tau \cdot g$ over the parameter range $\tau \in [0.5, 5.0]$,
$g \in [0.1, 1.0]$. The volume under this surface, divided by $N^2$, is the
universal-law prediction; the simulator estimates the same quantity by
uniform Monte Carlo sampling, providing an independent cross-check.}

\newcommand{\captionUTLpaneltwo}{%
\textbf{Queueing-theoretic specialisations.}
(a) M/M/1 case: the universal law (amber crosses) reproduces the classical
$TP^{-1}_{M/M/1} = \rho/(\mu(1-\rho))$ formula (blue line) across the full
range of utilisation $\rho \in [0,1)$, recovering the well-known divergence
as $\rho \to 1$.
(b) Jackson-network independence: composite throughput inverse via the
universal law (x-axis) versus the Jackson-network formula (y-axis); points
fall exactly on the identity line, confirming that for product-form networks
$TP^{-1}$ decomposes additively across stations.
(c) Cascade serialisation: as the cascade depth $k$ grows, the cascade-sum
formula and the universal law remain in agreement, validating that the
universal law subsumes the serial-composition rule
$TP^{-1}_\mathrm{cascade} = \sum_k TP^{-1}_k$.
(d) Surface plot of $\log_{10}(W)$ over the joint
$(\rho, \mu)$ parameter space; the steep ridge along $\rho \to 1$ visualises
the M/M/1 saturation regime predicted by the universal law.}

\newcommand{\captionUTLpanelthree}{%
\textbf{Five operating regimes.}
(a) Partition of the load coordinate $R \in [0,1]$ into the five regimes
\textsc{turbulent} ($R < 0.3$), \textsc{aperture-dominated} ($0.3 \le R < 0.5$),
\textsc{cascade} ($0.5 \le R < 0.8$), \textsc{coherent} ($0.8 \le R < 0.95$),
and \textsc{phase-locked} ($R \ge 0.95$). The boundary points $\{0.3, 0.5, 0.8, 0.95\}$
are inherited from the underlying Lagrangian formulation and validated
across heterogeneous domains.
(b) Histogram of regime occupancy under uniform sampling of $R$, confirming
the analytical width of each regime.
(c) Bifurcation curve $R(K)$ as the coupling strength $K$ varies; the
predicted (line) and measured (markers) curves agree to within numerical
tolerance, with the four regime boundaries marked by horizontal dotted lines.
(d) The \emph{regime cone}: a 3D surface of phase coherence over the
$(R, \epsilon)$ plane, where $\epsilon$ is an external perturbation amplitude.
The cone collapses toward zero coherence as $\epsilon$ grows, recovering the
classical perturbed-oscillator result.}

\newcommand{\captionUTLpanelfour}{%
\textbf{Critical slowing as a load indicator.}
(a) Log-log plot of measured relaxation time $\tau_\mathrm{relax}$ versus
distance $|R - R_b|$ from the nearest regime boundary $R_b$. The dashed
red line shows the universal scaling $\tau_\mathrm{relax} \propto |R - R_b|^{-1}$;
both measured (circles) and predicted (crosses) values follow this law over
several decades.
(b) Estimated distance-to-boundary $\hat d$ versus true distance $d$;
points coloured blue (correct) or red (incorrect) show the load-indicator's
classification accuracy. The estimator recovers the correct distance to
within tolerance for all trials.
(c) Per-trial relative error of the load indicator; the mean residual is
shown in the title.
(d) Three-dimensional surface of $\log_{10}(\tau_\mathrm{relax})$ over the
$(R, \epsilon)$ plane, demonstrating that the divergence of relaxation time
near regime boundaries is the structural signature exploited by the
load-indicator estimator.}

\newcommand{\captionUTLpanelfive}{%
\textbf{Cache extinction and federation composition.}
(a) Effective $TP^{-1}$ after cache deployment versus uncached $TP^{-1}$;
the dashed line marks the no-cache reference. The structural drop --- not
an asymptotic speedup --- corresponds to the cache extinguishing decision
classes whose lookups are constant-time. The mean speedup factor is shown
in the title.
(b) Federation composition: composite $TP^{-1}_\mathrm{fed}$ as the
number of federated kernels grows from $1$ to $20$. The curve saturates at
$\Sigma$ from below, confirming the multiplicative composition law
$1 - TP^{-1}_\mathrm{fed}/\Sigma = \prod_i (1 - TP^{-1}_i/\Sigma)$.
The maximum absolute residual against the closed-form formula is
recorded in the panel title.
(c) Marginal benefit of adding the $n$-th kernel; bars decrease
monotonically, demonstrating the diminishing-returns property
predicted by the federation law.
(d) Composition surface over $(n, TP^{-1}_\mathrm{per\;kernel})$ at fixed
$\Sigma = 100$; the upper plateau visualises the asymptotic ceiling
that no finite federation can exceed.}

\newcommand{\captionUTLpanelsix}{%
\textbf{Estimators and architecture invariance.}
(a) Phase-coherence estimator: empirical $R$ as a function of the input
concentration parameter $K$; the estimator recovers the expected
monotone increase from uniform ($R \approx 0$) to fully concentrated
($R \approx 1$).
(b) Lag estimator lower bound: measured $\min \tau$ versus the
theoretical lower bound $1/f_\mathrm{max}$; points lie above the dashed
identity line, confirming the kinematic constraint
$\tau \ge 1/f_\mathrm{max}$ across $30$ trials spanning seven orders of
magnitude in $f_\mathrm{max}$.
(c) Coupling estimator recovery correlation as a function of decision-class
count $N$; the dashed red line marks the $0.5$ acceptance threshold.
The mean recovery correlation is shown in the title.
(d) Cross-architecture invariance: predicted $TP^{-1}$ for four kernel
architectures (microkernel, monolithic, exokernel, unikernel) across
$5$ trials each. The universal law produces consistent predictions
under the architecture-specific normalisation, confirming that the
formula is architecture-agnostic up to choice of $\mathcal{N}$.}
